\chapter{Entrada/Salida Fundamental}
\label{chap:io}
{\huge$\bigstar$} Principiante

\section{La Función Format}
\begin{lispcode}
(format t "~A~%" "Hola")
(format nil "~A + ~A = ~A" 1 2 3)
(format nil "~D" 42)
(format nil "~B" 42)
(format nil "~X" 255)
(format nil "~5,2F" 3.14159)
(format nil "~@R" 42)
\end{lispcode}

\section{Print y Variantes}
\begin{lispcode}
(print "Hola") (princ "Hola") (prin1 "Hola") (terpri)
\end{lispcode}

\section{Lectura}
\begin{lispcode}
(let ((x (read))) (format t "Leí: ~A" x))
(read-line) (read-char) (read-from-string "(+ 1 2)")
\end{lispcode}

\section{Archivos}
\begin{lispcode}
(with-open-file (out "salida.txt" :direction :output :if-exists :supersede)
  (format out "Hola, archivo~%"))
(with-open-file (in "entrada.txt" :direction :input)
  (let ((c (read-line in nil nil))) c))
\end{lispcode}

\section{Redirección de Streams}
\begin{lispcode}
(with-output-to-string (out) (format out "Esto va al string"))
(with-input-from-string (in "(1 2 3)") (read in))
\end{lispcode}

\begin{examplebox}[title=Formato tabular]
\begin{lispcode}
(defun tabla-cuadrados (&optional (n 10))
  (format t "~&~@A~20@A~30@A~%" "N" "N^2" "N^3")
  (dotimes (i n)
    (let ((x (1+ i)))
      (format t "~&~D~20@D~30@D~%" x (* x x) (* x x x)))))
\end{lispcode}
\end{examplebox}

